% !TEX root = ../thesis_book_0421052099_appendices.tex
\chapter{Implementation Notes}\label{app:implementation}

% The four techniques Section~\ref{sec:implementation} points here for. Written
% as reusable techniques rather than as a changelog: each one is a practice
% another project could adopt, and the bug that motivated it is the evidence.

\Cref{sec:implementation} names three practices that carry more reproducibility
weight than the dependency list, and promises they are recorded here as
techniques rather than as incidents. Those are the first three sections below.
Three further techniques follow, which govern how the agent itself is built and
which the earlier three depend on. In each case the defect that forced the
practice is given as evidence rather than as narrative.

\section{Records Self-Describe}\label{app:self-describing}

Every record \textsf{RunLogger} writes carries, alongside its result, the
backbone description, a hash of the budget configuration, the run configuration,
and the source commit. On the main matrix, the ablations and the smoke baselines
--- $6{,}060$ records --- the run configuration also names the system, because
the runner adds that key to it rather than \textsf{RunConfig} carrying it. The
development sweep of \Cref{sec:dev-tuning} goes through the same logger and
carries the same four stamps. But nothing there adds that key. So those $900$
records identify their system by file name alone. The backbone qualification of
\Cref{sec:backbone} predates the logger and writes its own records. Those
records carry the backbone description and none of the rest.

The cost is a few hundred bytes per question. The benefit is that provenance
stops depending on where a file sits or what it is called: a record found in the
wrong directory, or copied out of one, still states what produced it. The
cache-replay anomaly of \Cref{sec:backbone} is both what forced the practice and
what shows its limit --- the question was which model had produced a set of
records, and the backbone stamp answered it. But those same records carry
neither the configuration nor the commit. That gap is the one the logger was
written to close.

The generalisation is that a file path is not metadata. Directory layout is
edited, archived and reorganised over a project's life. Every reorganisation
is an opportunity to detach a result from the conditions that produced it.
Anything a reader must know in order to interpret a record belongs inside the
record.

\section{Frozen Artifacts Are Committed; Regenerable Ones Are
Not}\label{app:frozen-vs-regenerable}

The distinction is not about file size. It is about whether the artifact can be
recreated.

The certified question sets, their half-splits, and every result file embody
spent money and runs that cannot be repeated identically. They are committed.
The graph import, the embedding matrices and the triple index are deterministic
products of committed scripts run over public data, and are not committed even
though rebuilding them takes hours.

Applying the rule requires honesty about the second category. An artifact is
only regenerable if the script that generates it is committed \emph{and} its
inputs are still obtainable. Where either fails, the artifact is frozen
regardless of how it was made. \Cref{sec:implementation} records the one
case in this project where that judgement had to be made after the fact: a
generated summary in the development-set archive that predates a fix applied to
the runs it summarises.

\section{Cache-Identity Verification}\label{app:cache-identity}

This is the technique of most methodological interest, and the one most readily
reusable outside this thesis.

The problem it solves is stated easily and answered badly by ordinary testing. A
refactor is made to code on the frozen experimental path. Tests pass. Have the
reported numbers changed? Re-running the experiment answers the question. But a
re-run that produces the same numbers is weak evidence, because it is exactly
what a re-run would produce if the change had also altered which questions
happened to succeed.

The stronger check exploits the response cache. Every model call is keyed by a
hash over the model identifier, the sampling parameters, and the full prompt
text. A code change that does not perturb the frozen path therefore issues
\emph{byte-identical prompts}. That means every call must be a cache hit. So:

\begin{quote}
Re-run the affected questions and assert that the number of cache hits equals
the number of model calls, and that both are non-zero.
\end{quote}

A single cache miss proves that some prompt differed, and therefore that the
change was not behaviour-preserving --- before any output is compared. The
assertion is stronger than output comparison because it tests the input to the
non-deterministic component rather than the output of it, and it is cheap
because a fully cached re-run costs nothing.

For this to work the cache hit must be indistinguishable from a live call inside
the agent. The budget check still runs, the call counter still increments, and
the original token counts are replayed into the meter from the stored record. A
cached re-run therefore reproduces the original budget snapshots exactly rather
than reporting zeros. That fidelity is what makes the identity assertion
meaningful. A separate counter records hits. So live and replayed runs remain
distinguishable in analysis.

Anyone lifting this design should carry one defect with it. The replay restores
prompt and completion tokens but not \textsf{reasoning\_tokens}: the cache
stores that field and never adds it back to the meter. The backbone here runs at
reasoning effort \textsf{none}, and the field is zero in every one of the
$6{,}960$ agent runs committed for this work. So ``exactly'' is literally true
of these runs. Under any reasoning-effort setting, the replayed snapshot would
understate the budget, and the identity assertion would be checking a weaker
claim than it appears to. \Cref{sec:pre-registration} gives the reason the
defect is recorded rather than repaired.

\Cref{sec:pre-registration} reports the one occasion this certificate was
demanded and returned a clean result, for the \textsf{use\_planner} flag that
must be a no-op when disabled. \Cref{sec:instrumentation} specifies the cache
whose behaviour makes the certificate meaningful.

\section{Config Injection and the Partial-Edit
Hazard}\label{app:config-injection}

The agent is a graph of nodes whose signatures the framework, not the author,
controls. Passing a run configuration into a node therefore means either closing
over it at construction time or threading it through the framework's own
injection mechanism. The first is easier and is a trap: a node that closes over
a configuration value captures whatever was in scope when the graph was built.
For a sweep, that means the first condition's value silently serving every
condition.

The practice that follows is narrow and worth stating plainly. \emph{A
configuration value must reach the code that uses it by exactly one route.}
Where a value can arrive by two routes, an edit that updates one of them leaves
the other in place. The resulting run is neither the old condition nor the
new one. This is the partial-edit hazard, and it is not caught by tests, because
each route works in isolation.

The defence used here is an assertion at the entry of every node that receives a
configuration, checking the type of what actually arrived. It converts a silent
wrong-condition run into a crash on the first question. A run that dies
immediately costs minutes; a run that completes under the wrong condition
contaminates a table.

\section{Routers Read, Nodes Write}\label{app:state-discipline}

The state object is mutated by nodes and inspected by routers. Allowing routers
to write as well produces control flow that cannot be reconstructed from the
trace, because the state a node observes was modified by the routing decision
that reached it.

The discipline is one sentence: \emph{routers may read state and must not write
it; nodes may write state and must not route.} It costs a small amount of
duplication --- a value a router needs must be written by the node before it,
under a name the router agrees on --- and buys the property that the trace is a
complete account of the run. Every claim in \Cref{sec:erroranalysis} about where
a trajectory went wrong depends on that being true, because the census reads
traces, not code.

Two categories of scratch key exist in the state for this reason. Node outputs
are the record. Router inputs, prefixed with an underscore, exist so that a
router can decide without recomputing what a node already knew.

\section{Instrumentation Decides What Can Be
Concluded}\label{app:instrumentation-gap}

The verifier persists the claims it rejects. It does not persist the claims it
accepts.

That decision was made early, on the reasonable ground that rejections are the
interesting events: they are what the layer does, and they are few. Its
consequence surfaced only at analysis time. Wrongful rejection --- a true claim
hedged --- is recoverable from the logs at census scale, because every rejection
is on record. Wrongful acceptance --- a false claim passed --- is not, because
an accepted claim leaves no trace distinguishable from a claim that was never
tested.

The thesis therefore reports one polarity of verifier error as a rate and the
other from a single specimen recovered by manual trace reconstruction. Since
wrongful acceptance is the polarity that bears directly on a hallucination
claim, this is the limitation named first in \Cref{sec:limitations-final}.

The technique to extract from this is not a fix. The fix is obvious and is the
first item of \Cref{sec:future}. It is the prior question:

\begin{quote}
Before freezing an instrument, ask which claim it will be asked to support, and
whether the records it keeps can support the negation of that claim as well as
its assertion.
\end{quote}

Logging the rejections was sufficient to describe what the layer does. It was
not sufficient to bound what the layer misses. Only the second supports a
claim about hallucination. The cost of logging both was negligible; the cost of
not doing so was a limitation that no amount of later analysis could remove.


\section{The Smoke-Set Validation Record}\label{app:smoke}

Two tuning sets appear in this thesis. They are not the same. This section
uses the \emph{smoke set}: twenty questions drawn from the training and
validation splits, never from test, used once to validate the assembled
framework's mechanics. The $80$-question \emph{development set} of
\Cref{sec:dev-construction}, on which the hyperparameters were swept, is a
separate and larger construction. Only it carries any reported number. The
smoke set was built to exercise specific machinery rather than to sample the
benchmark distribution: six single-hop questions, eight two-hop compositions,
three conjunctions, two whose shortest path traverses a mediator, and one about
a deliberately non-existent entity whose correct answer is a refusal.

The first end-to-end run answered roughly three to four questions correctly; the
final configuration answers thirteen to fourteen. That interval is not
incidental tuning --- it is where several of the mechanisms described above were
shown to be necessary. It is also the one place in this thesis whose
\emph{evidence} is not fully archived. What survives is one record file,
\textsf{results/phase3/smoke20\_a0.5\_t0.2.jsonl}, scoring thirteen correct, six
hedged and one wrong at $\alpha = 0.5$ rather than the frozen $0.7$. What does
not is the before-state, read from a console log and never committed. So the
three-to-four figure and the root causes of \Cref{tab:rootcauses} rest on that
reading alone. Neither number enters any result. That is why the omission was
tolerable and why the figures are given as ranges.

The first run produced no crashes, terminated cleanly on every question, and
abstained rather than confabulated in almost every failure --- which was the
actual exit criterion for the framework's mechanics. Its sixteen or seventeen
failures collapsed into five systematic causes, summarised in
\Cref{tab:rootcauses}.

\begin{table}[!tb]
  \begin{center}
    \caption[Root causes of failure in the first end-to-end run]{Root causes of
    failure in the first end-to-end run, and the mechanism each one motivated.
    Every fix is described in the section indicated.}
    \label{tab:rootcauses}
    \footnotesize
    \begin{tabular}{|L{52mm}|L{52mm}|l|}
      \hline
      \textbf{Observed failure} & \textbf{Mechanism introduced} &
      \textbf{Section} \\
      \hline
      Backtracking was a deterministic no-op: the restored frontier produced
      identical scores, facts, and decisions until the budget drained &
      Ban list on abandoned \textsf{(anchor, relation, direction)} triples &
      \Cref{sec:backtracking} \\
      \hline
      The evaluator's fact window showed the most recent thirty triples, so a
      large successful expansion pushed the answer out of view &
      Relevance-ranked fact window & \Cref{sec:evaluator} \\
      \hline
      The relation list, capped at 80 by descending fanout, hid low-fanout
      precise relations behind hub relations &
      Cap raised to 300, plus the administrative-predicate blocklist &
      \Cref{sec:tool-api} \\
      \hline
      Entities resolved while an objective was still incomplete were discarded &
      \textsf{candidate\_answers} accumulated on every evaluation &
      \Cref{sec:evaluator} \\
      \hline
      The frontier was truncated in insertion order, so the agent drifted away
      from the anchor that mattered & Frontier sorted by edge score before
      truncation & \Cref{sec:explorer} \\
      \hline
    \end{tabular}
  \end{center}
\end{table}

Three observations from this exercise bear directly on later chapters.

\textbf{The two confabulations in the first run were caused by a retrieval
defect, not a generation defect.} Two questions about national currencies
produced a confident ``United States Dollar'': the relation cap had hidden the
\textsf{currency\_used} relation, the beam surfaced triples about GNI per capita
denominated in dollars, and the answerer pattern-matched. No retrieved triple
supported the answer. After the cap was raised both questions resolved
correctly, at roughly a third of the token cost. These two traces are the
motivating example for \Cref{sec:verification-layer}: a claim-level check
against the traversed triples would have caught them regardless of the retrieval
defect. That is precisely the argument for verifying before emitting.

\textbf{The groundedness filter on resolutions was validated by a failure.} On
one question the evaluator listed nine correct countries, none of which appeared
in any retrieved triple, because the beam had gone down an unrelated branch. The
filter rejected all nine, and the system abstained. That cost a point. The
rejection also demonstrated, one layer earlier than intended, that a language
model inside an agentic scaffold will assert entities it never retrieved, and
that catching this requires an explicit check --- nothing about the model's
confidence distinguishes the nine countries it recalled from entities it
actually found.

\textbf{One instrumentation defect was found and removed.} A counter reporting
how many banned expansions had been skipped was computing the number of anchors
instead, reporting a plausible-looking number unrelated to what it claimed to
measure. It was corrected rather than deleted, and it motivates a standing rule
adopted for the remainder of the work: a metric that has never been checked
against a case where its correct value is known independently should not be
trusted, and should not be analysed.
\subsubsection{The Complete Navigation
Algorithm}\label{sec:navigation-algorithm}

\Cref{alg:agr} states the loop.

\begin{algorithm}[!tb]
  \caption[AGR navigation]{AGR navigation (verification layer abstracted; see
    \Cref{sec:verification-layer})}
  \label{alg:agr}
  \begin{algorithmic}[1]
    \REQUIRE question $q$, graph $G$, budgets $B$, blend $\alpha$, threshold
      $\tau$
    \ENSURE answer $A$ with supporting triples $S$
    \STATE $P \gets \textsc{Plan}(q)$; \textbf{if} invalid \textbf{then}
      $P \gets [\,q\,]$
    \STATE $F \gets \textsc{SeedAnchors}(P)$;\quad
      $T \gets \emptyset$;\quad $C \gets \emptyset$;\quad
      $\mathit{stack} \gets \emptyset$;\quad $\mathit{ban} \gets \emptyset$
    \WHILE{budgets permit}
      \STATE push $(F, \mathit{depth}, \sigma)$ onto $\mathit{stack}$
        \COMMENT{$\sigma$ from the previous pass; $0$ initially}
      \STATE $R \gets \bigcup_{e \in F} \textsc{GetRelations}(e)$
      \STATE $\mathit{scored} \gets \alpha\cos(\cdot) + (1-\alpha)\mathrm{LLM}
        (\cdot)$ over $R$ against the active objective
      \STATE $E \gets$ top-$\beta$ of $\mathit{scored}$ per anchor, excluding
        $\mathit{ban}$
      \STATE $T \gets T \cup \textsc{GetNeighbors}(E)$;\quad
        $F \gets$ top-$m$ new entities by score, \textbf{or} $F$ unchanged if
        that set is empty \COMMENT{the dead-end trigger fires on the same
        condition and routes away before $F$ is reused}
      \STATE $\sigma \gets \max(\max_{c \in E}\mathit{scored}(c),\;
        \max_{c \in R}\cos(\cdot),\; \max_{c \in R_K}\mathrm{LLM}(\cdot))$
        \COMMENT{signal maximum}
      \STATE $(d, \mathit{done}, \rho) \gets \textsc{Evaluate}
        (\text{objective}, \textsc{TopFacts}(T))$
      \STATE $\rho \gets \{x \in \rho : x \text{ appears in } T\}$;\quad
        $C \gets C \cup \rho$
        \COMMENT{groundedness filter}
      \IF{$\mathit{done}$}
        \STATE advance $P$; \textbf{if} no objective remains \textbf{then}
          $d \gets \textsf{answer}$ \textbf{else} $F \gets \rho$
      \ENDIF
      \IF{$|T_{\mathrm{new}}| = 0$ \OR $\sigma < \tau$ \OR
          $d = \textsf{backtrack}$}
        \STATE $\mathit{ban} \gets \mathit{ban} \cup E$;\quad
          $(F, \mathit{depth}) \gets$ pop highest-scoring snapshot
      \ELSIF{$d = \textsf{answer}$}
        \STATE \textbf{break}
      \ENDIF
    \ENDWHILE
    \STATE $(A, S) \gets \textsc{VerifyAndAnswer}(q, T, C)$
      \COMMENT{\Cref{sec:verification-layer}}
    \RETURN $(A, S)$
  \end{algorithmic}
\end{algorithm}


\section{Failure Modes of the Verification Layer}\label{app:failure-modes}

This section is analytical: it enumerates the ways the layer can be wrong that
follow from its design, independently of how often they occur.
\Cref{sec:verifier-errors} reports which of them actually fired. In the
vocabulary of \Cref{sec:polarities}, a claim is \emph{wrongly rejected} when the
layer withholds support the environment would in fact provide. It is
\emph{wrongly accepted} when it certifies a claim the environment does not
support.

\subsection*{Mechanisms of Wrongful Rejection}

\textbf{Composite claims.} The decomposer emits claims corresponding to what the
draft asserts, not to every relationship the draft relies on. When a draft
asserts a filtered set --- ``these teams were founded before 1996'' --- the
decomposition is about the founding. The membership relation that makes each
team a candidate may never be claimed at all. Every team's only claim is then
the unverifiable one, the filter drops every entity, and a draft containing four
correct answers becomes a hedge. This is the Harbaugh question of
\Cref{tab:census}: \emph{a composite assertion whose unverifiable component
drags down its verifiable component}. The conservative rule is defensible for a
hallucination-mitigation system, and it is the layer's most expensive single
behaviour on the development set.

\textbf{Endpoints outside the traversal.} $\mu$ is built from traversed triples
only. So a claim about an entity the agent never walked to skips both structural
routes regardless of what the graph contains, and is judged solely by a model
against a fact window that by construction does not mention it.

\textbf{A question-ranked window for a claim-level judgement.} The entailment
check's evidence is ranked for relevance to the question, not to the claim under
test. So peripheral claims are systematically disadvantaged.

\textbf{Conservative failure.} Any exception in the entailment call marks all
pending claims unsupported: infrastructure failure is recorded as rejection.

\textbf{Name collisions in $\mu$.} The mapping keeps the first identifier seen
for a name. So two distinct graph entities sharing a surface form collapse to
one, and adjacency is tested for the wrong node.

\subsection*{Mechanisms of Wrongful Acceptance}

\textbf{Relation-agnostic adjacency.} Both structural routes ignore $r$ and
ignore direction. So any edge between the endpoints certifies any claimed
relation between them: a claim that X is Y's mother is accepted on an edge
recording that X is Y's child, and a claim that a film was directed by a person
is accepted on an acting credit. That symmetry is the direct cost of allowing
free-text relations and it is the layer's principal acceptance risk. Its scale
is larger than \Cref{sec:structural-check}'s development-set figure invites,
because neither route matches on the relation: the population exposed is every
claim the two of them accept between them. On test, that is somewhere in $[39,
2{,}008]$ of the $2{,}008$ accepted claims. The lower end is what the log pins
down; the upper end is what the mechanism permits. Nothing measured in this
thesis narrows the interval. \Cref{sec:future-logging} is the change that
would.

\textbf{The mediator hop.} \textsf{verify\_connection} accepts a path through
one mediator node. Because a Freebase mediator encodes a single $n$-ary fact,
the endpoints are genuinely co-participants in one fact. But which role each
plays is exactly what the mediator encodes and the test discards: an actor, a
film, a character, and a period are mutually ``connected'' through one
performance node, and any claim relating any two of them is certified.

\textbf{Vacuous certification of empty decompositions.} A draft with no claims
has no unsupported claims. So it is routed \textsf{grounded}. On the development
set this describes $10$ of $80$ questions. The label is not false --- an answer
that asserts nothing asserts nothing unsupported. But it must not be read as
a positive verification. Any statistic computed over \textsf{grounded}
outcomes inherits this.

\textbf{Entities that appear in no claim are never checked.} This is the
broadest mechanism listed here. Verification operates on claims, while the
answer's entity list is a separate output of the same call. So an entity the
drafter named but the decomposer never wrote into a claim passes through
untouched on the grounded path. It is deliberately retained by the permissive
branch of the filter on the repair path (\Cref{sec:entity-filtering}). Since
$77$ of $80$ development questions took the grounded path, the layer's guarantee
over emitted entities is in practice mediated entirely by decomposition recall.
That recall is itself a language-model output subject to no check.

\textbf{Grounded-but-vacuous answers.} This is the boundary of what triple-level
verification can do. An answer whose every claim is supported can still fail to
answer the question: ``the senators from New Jersey are the United States Senate
members associated with New Jersey'' decomposes into claims the graph genuinely
supports, and is certified. No claim-level check can catch this, because the
defect is not in any claim but in the relationship between the claims and the
question. The anti-vacuity provision of \Cref{sec:claim-decomposition} addresses
it at drafting, which is mitigation rather than verification.

\textbf{Structural acceptance is final.} A claim accepted by adjacency is never
examined semantically. The entailment check, the only component that looks at
meaning, sees only what the structural routes could not resolve.

\subsection*{What the Layer Is}

Stated at the level of precision the preceding list requires: the layer
establishes that the entity pairs a draft's claims relate are connected in the
environment, and on the repair path it makes the emitted entity list a
deterministic function of those verdicts rather than a second generation. It
does not establish that the claimed relation is the one the edge records, that
an entity absent from every claim was grounded at all, that the answer is
responsive to the question, or that a claim it could not check is false. Those
boundaries are properties of the design rather than defects discovered in it.
\Cref{sec:erroranalysis} reports how much of the residual error they account
for.


\section{Backbone Qualification: Four Determinism Rounds}\label{app:backbone}

Selection between \textsf{gpt-4.1-mini-2025-04-14} and
\textsf{gpt-5.4-mini-2026-03-17} was made by a qualification script that runs
both candidates over the same twenty development questions and reports tokens,
calls, latency, reasoning tokens, errors, and a determinism probe. The probe
re-runs a subset of questions and counts how often the \emph{decision signature}
--- a hash of the plan, the expanded relations, the evaluator verdicts, and the
answer, not of timings or scores --- repeats exactly. The first round reported
$3/3$ matches for 4.1-mini against $1/3$ for 5.4-mini; a second round with the
cache disabled reversed it, $2/3$ against $3/3$. The reversal, not either
number, is the finding.

Temperature-zero decoding on a hosted API is greedy, not deterministic: the
logits are not bit-stable across calls, because floating-point reduction order
depends on how a request is batched with other traffic. A wobble of order
$10^{-6}$ flips the argmax when the top two tokens are near-tied. A sequential
agent then \emph{amplifies} that, since one divergent token in a planner call
changes a sub-objective's wording, which changes its embedding, which changes
the beam. A model that is $98\%$ stable per call can be $70$--$90\%$ stable per
trajectory. With three probes per round, an estimator of a rate in that band
bounces between $1/3$ and $3/3$ by sampling variance alone. A tiebreaker was
therefore pre-registered before rounds three and four: \emph{if the two
candidates finish within three matches of each other, take
\textsf{gpt-5.4-mini}} on longevity grounds, the 4.1 line being
deprecation-adjacent. Both rounds returned $2/3$ each, the pooled score was
$9/12$ against $8/12$ --- one match apart, inside the band. So the rule fired,
and the backbone was frozen on longevity.


\section{Output-Contract Instrumentation Defects}\label{app:contract-defects}

Three instrumentation defects in this area are recorded rather than quietly
fixed, in keeping with the same standing rule. \emph{First}, a Python
falsy-empty-list trap: the answerer substituted the entire traversed set for the
supporting triples whenever the verifier's list was empty, treating ``the
verifier ran and matched nothing'' identically to ``the verifier never ran''. So
a development trace recorded $1{,}927$ supporting triples for an answer that
asserted nothing. That trace was read before the fix and never committed. So the
figure cannot be recomputed here. The largest support count in any committed
record is $133$, and neither enters a result. The fix tests for \textsf{None}
rather than for falsity. Its stated rationale was wider than the code, though:
the intended three-way distinction between an absent key and an empty list is
unreachable, because the state initialiser seeds \textsf{supporting\_triples} to
an empty list. So the fallback branch is dead code guarding a state the
initialiser rules out.

\emph{Second}, and unrepaired, a counter in the verifier's trace named
\textsf{n\_structural} counts every supported claim including those certified by
entailment. So its name misdescribes it. Under the standing rule it is fenced:
no table, figure or reported result derives from it. The route-by-route
counts of \Cref{sec:structural-check} come from the tool invocation log instead.
The fence has one exception, named here because a rule applied silently is
indistinguishable from a rule applied selectively --- the third defect reads the
counter once, to bound how often a duplicate could arise. That is admissible
only because the direction of its error is known.

\emph{Third}, a mismatch between \Cref*{alg:verify} and the code it describes.
The algorithm accumulated $E$ by set union; the implementation appends. Two
claims in one draft sharing an endpoint pair --- ``X was born in Y'' and ``X
died in Y'', decomposed separately but resolving to the same two identifiers ---
therefore each append the same matching triples. So
\textsf{n\_supporting\_triples} counts matches rather than distinct triples. The
algorithm is corrected here to say what the code does, since the code is what
produced the results. How often it bites cannot be recovered, the claim
endpoints not being persisted either. But the opportunity can be bounded: a
duplicate needs two claims matched by the adjacency route, and over the $908$
verifier firings of the development and test sets, $378$ report two or more
supported claims. That $378$ is an over-count of the exposure and knowably so,
since the counter includes claims certified by the other two routes, neither of
which appends anything. So \textsf{n\_supporting\_triples} is an \emph{upper}
bound on distinct supporting triples. The median of $3$ and mean of $4.1$
inherit that.


\section{Where Each Budget Is Enforced}\label{app:budget-sites}

Enforcement is confined to routers and the client. That is what keeps it
auditable. But the mapping is not one site per budget. The model-call budget is
checked inside the client before every call. So no node can bypass it.
Exceeding it raises an exception that each node catches into a degraded path
rather than a crash. The backtrack budget is checked in the post-evaluation
router alone. Beam width and the anchor cap are not checks but slices: beam
width in the explorer, the anchor cap in the explorer and again in the
evaluator, where resolved entities become the next frontier. The remaining three
are each checked in two places. In every case the second site exists because
the first does not cover the verify--explore cycle. \emph{Depth} and
\emph{wall-clock} are checked in the post-evaluation router and again in the
post-verification router, since the retry edge re-enters the explorer without
passing through the first. Without the second, a repair would silently buy an
expansion past the depth budget --- \Cref{sec:repair} records the development
trace that exposed exactly that. \emph{Verification iterations} are checked in
the post-verification router and again inside the verifier, which will not
append a repair objective it has no budget to pursue. The two sites read the
same counter through the same predicate. So neither can drift from the other.


\section{Backtracking: Recorded Deviations}\label{app:backtrack-deviations}

\textbf{The implementation departs from what this section specifies in three
places, and all three are recorded rather than repaired}, under the standing
rule of \Cref{sec:design-validation}, since repairing them would change the
frozen results this thesis reports. \emph{First, $\sigma = 0$ is read as ``no
signal recorded''.} A missing signal defaults to a value above any usable
$\tau$. That is right when no expansion has run yet. But $0.0$ is falsy in
Python. So a genuine signal maximum of exactly zero takes the same branch, and
the trigger does not fire at the one point where the formula says it most
clearly should. This is the same falsy-empty trap \Cref{sec:output-contract}
records in the answerer, found in a second place. \emph{Second, at $\alpha \geq
1$ only the first term is computed}, the scorer taking an early return when the
model term is disabled. So $\sigma$ collapses to the blended maximum over
expanded edges. The same early return is taken when the candidate list is empty.
There the two auxiliary maxima are not merely unrecorded but \emph{stale},
left holding whatever the last pass wrote. Since the runtime memoises one
scorer instance per $\alpha$ for the process, those values persist across every
question in the run. Both defects have narrow practical reach, because a signal
maximum of exactly zero and an empty candidate list both trip the \textbf{dead
end} trigger, which is evaluated first. The masking is the order of the
triggers rather than anything about the values. That is a weak reason for the
behaviour to be correct. The second one does have a visible effect at $\alpha =
1$: it measures $\sigma$ over a \emph{smaller set}, the beam-selected
non-banned edges rather than every candidate the anchors offered. So after a
backtrack has banned the best relation, $\sigma$ falls where it would have held
steady in the reference. That is exactly the embedding-only ablation.
\Cref{sec:abl-embonly} reports what it does to that condition's backtrack count
instead of attributing the difference wholly to $\tau$.

\emph{Third, the ban list is narrower than the jump it has to cover}, and this
is the deviation with a measurable consequence. What the backtracker bans is the
\emph{most recent} explorer pass; what it restores is the \emph{highest-scoring}
snapshot, frequently older than that pass. Everything expanded in between stays
unbanned. So on exactly those backtracks the explorer can be handed back a
frontier with all of its previously selected edges still available, and
re-selects them. Over the main matrix, $53$ explorer passes --- $15$ on WebQSP
and $38$ on CWQ --- expanded an \textsf{(anchor, relation)} set identical to one
the same question had already expanded: exactly the repetition the ban list
exists to prevent. The misalignment that produces it is visible in at least $65$
of the $262$ backtracks, where the restored depth is not the depth the most
recent snapshot carried.\footnote{A lower bound rather than a count. Identifying
the popped snapshot exactly would need the scores. The trace does not
carry them. Depth it does carry. The depth series reconstructs exactly against
the budget snapshot on all $800$ records. A pop that restores an older snapshot
sitting at the same depth is indistinguishable from a most-recent pop and is
counted as one here.} The unit test that should have caught this compares the
first expansion set against the second. So it is structurally incapable of
observing a gap that opens only when the stack is popped out of order: the test
passes, and what it establishes is narrower than its name suggests.


\section{The Development Sweep: Curve Shape and Threshold}\label{app:sweep}

Three readings, of which only the first is the one the sweep was run to obtain.

\textbf{$\alpha = 0.7$ is the maximum. The curve is not clean.} Hits and F1
rise from $\alpha = 0.3$ to $\alpha = 0.7$ and fall sharply at $\alpha = 1.0$.
That is the expected shape: embedding similarity alone ($\alpha = 1.0$) is a
weaker guide than the blend. Weighting the language-model judgement too
heavily ($\alpha = 0.3$) makes the score noisy. But the wrong-answer column does
not follow the hits column. It reads $7, 8, 5, 6$ across the four $\alpha$
values, peaking at $\alpha = 0.5$ --- the condition that also has the fewest
hedges of the verified rows. At $\alpha = 0.5$ the system commits more often and
is wrong more often. That irregularity is reported rather than smoothed
away, because a presented curve that hides a known kink is exactly what a
defence finds. It is one sample of $80$ questions. The difference is three
questions wide. It is noted, not interpreted.

\textbf{$\tau$ does almost nothing on this set. That is consistent with its
design.} At $\alpha = 0.7$ and $\alpha = 0.5$ the two $\tau$ values produce
identical hits, hedges, and wrong counts. At $\alpha = 0.3$ they differ only in
backtrack count ($16$ against $18$). At $\alpha = 1.0$, one wrong answer becomes
a hedge. Backtracks rise from $22$ to $26$. This is what the signal-maximum
formulation of \Cref{sec:backtracking} predicts: the trigger fires only when the
blended score, the embedding maximum, \emph{and} the language-model maximum are
all below threshold. That conjunction is rare when the language-model term is
contributing. The $\alpha = 1.0$ row is where it has the most to do. It is
exactly there that its effect is visible. The value $\tau = 0.20$ was frozen on
the basis of a sweep in which it changed almost nothing. That is stated plainly
rather than presented as a tuned choice.

\textbf{Verification costs accuracy on this set.} Against the frozen condition,
the draft-only row scores $66$ hits to $64$ and F1 $0.797$ to $0.769$. After the
answer-entity filter of \Cref{sec:entity-filtering} was corrected --- a change
made on development data before any test question ran --- the frozen condition
scores $65$ hits, $11$ hedges, $4$ wrong, F1 $0.785$. Across those two rows the
layer costs one correct answer and removes no wrong ones on this set of $80$.
That comparison spans the correction, since the draft-only condition was never
re-run after it, and is therefore the closest available reading rather than a
matched pair. \Cref{sec:ablations} finds the same thing at $n \approx 200$ on
test data. \Cref{sec:findings} explains why that is the expected result
rather than a disappointing one.

One reproducibility note belongs with this table. Only the frozen condition was
re-run after the filter correction. So \Cref*{tab:sweep} is the pre-correction
sweep throughout --- internally consistent, since all eight conditions ran under
the same code. The committed summary file for that directory also predates the
correction. So rescoring the archived runs yields $65/11/4$ for the frozen
condition where the summary records $64/11/5$. The archived per-question records
are authoritative.


\section{Static GraphRAG: The Fanout Cap and Its Population}\label{app:fanout}

\textbf{The radius is not the only thing that bounds this baseline. The fanout
is, and it binds harder.} Expansion retrieves at most $100$ edges per topic
entity and at most $20$ per mediator hop. The Cypher that fetches them
carries no \textsf{ORDER BY}. So on any entity with more than $100$
post-blocklist edges, the $100$ retrieved are an arbitrary sample of the
neighbourhood rather than its best or its first. The population the cap acts on
is exactly identifiable, because this baseline seeds from the dataset's
annotated topic entities and expands nothing else: the $1{,}032$ topic mentions
across the $800$ test questions resolve to $711$ distinct entities, each
matching exactly one node at the resolver's exact-match stage
(\Cref{sec:entity-resolution}). So the full system necessarily seeded the same
nodes, and the tool log records which. Their median post-blocklist degree is
$131$, the ninetieth percentile $1{,}623$, the largest $136{,}742$. And $55.1\%$
carry more edges than the cap retains. Since a question carries $1.29$ topic
entities on average, the question-level share is measured separately: on
$72.5\%$ of the $800$ questions at least one topic entity is truncated, and on
$59.0\%$ every one of them is. For a hub topic the baseline answers from a small
fraction of the available neighbourhood --- on the largest, well under one
percent --- chosen without a stated criterion. The consequence for the hub-noise
diagnosis of \Cref{sec:hop-results} is that the answer may never have entered
the sample the reranker saw. Because the ordering is unspecified rather than
seeded, this is the one place in the harness where a re-run is not guaranteed to
reproduce the same retrieved set, though it did across the runs recorded here.

In fairness the same property holds of the full system's own
\textsf{get\_neighbors}, whose Cypher also carries \textsf{LIMIT} without
\textsf{ORDER BY} (\Cref{app:tool-neighbors}). The two are not comparably
exposed: the agent's cap is $200$ rather than $100$ and is applied per relation
rather than per entity. So it binds on $3.3\%$ of its $16{,}301$ neighbour calls
against the baseline's $55.1\%$ of topic entities. \Cref{sec:tool-api}'s
description of the tools as parameterised Cypher templates should be read with
that qualification: the parameters are fixed, the row order under a binding cap
is not.


\section{The Entity-Filter Attribution Census}\label{app:filter-census}

\begin{table}[!tb]
  \begin{center}
    \caption[The complete attribution census]{The complete attribution census:
    every development-set question whose answer differed between claim checking
    and its ablation, at $\alpha = 0.7$, $\tau = 0.20$. Three questions, one run
    per condition. ``Pre-fix'' is the original design in which the rewrite call
    regenerated the entity list. ``Post-fix'' is deterministic filtering. Per
    \Cref{sec:hedge-accounting}, the three fabricated-entity questions --- the
    \textsf{unanswerable\_fake} bucket, whose correct answer is a refusal ---
    hedge in every condition and are counted as hedges rather than excluded, so
    the totals are over all $80$. These are not the unanswerable-in-environment
    class of \Cref{sec:unanswerable}, which is a separate bucket of four whose
    questions do have true answers.}
    \label{tab:census}
    \footnotesize
    \begin{tabular}{|L{29mm}|L{35mm}|L{35mm}|L{29mm}|}
      \hline
      \textbf{Question} & \textbf{Draft-only} & \textbf{Verified, pre-fix} &
      \textbf{Verified, post-fix} \\
      \hline
      Harbaugh's teams founded before 1996
      \newline \emph{gold:} 4 teams &
      4 gold teams \emph{plus} Baltimore Ravens; scores as a hit &
      \textsf{[\,]} --- hedge &
      \textsf{[\,]} --- hedge \\
      \hline
      Party of Hitler and Carl Voss
      \newline \emph{gold:} Nazi Party &
      \textsf{Nazi Party} --- hit &
      \textsf{Adolf Hitler}, \textsf{Nazi Party}; hit, but asserts the
      question's subject &
      \textsf{Nazi Party} --- hit \\
      \hline
      Franco's university, founded 1701
      \newline \emph{gold:} University Yale &
      \textsf{University Yale} --- hit &
      \textsf{Yale University} --- \textbf{wrong} &
      \textsf{University Yale} --- hit \\
      \hline
      \multicolumn{4}{|l|}{\textbf{Whole run} --- all 80 questions; the three
        counts partition the set} \\
      \hline
      hits / hedges / wrong & 66 / 10 / 4 & 64 / 11 / 5 & 65 / 11 / 4 \\
      \hline
    \end{tabular}
  \end{center}
\end{table}

The classification is unambiguous. \emph{No entailment verdict was wrong}. No
repair exploration was involved --- as \Cref{sec:repair} records, none ran.
All three differences trace to one component: the rewrite call, failing three
different ways. On Harbaugh it over-deleted, discarding four grounded-correct
team names along with the unsupported one. On Hitler/Voss it laundered the
question's subject into the answer list --- an entity the draft never asserted.
On Franco it silently respelled the graph's \textsf{University Yale} as the
parametric \textsf{Yale University}, converting a defensible qualified answer
into a scored error.

The finding is worth stating in one sentence, because it is the conclusion this
chapter's design rests on:

\begin{quote}
The verification logic was sound; its repair path free-generates entities from a
language model, which reintroduces exactly the ungroundedness the layer exists
to prevent.
\end{quote}

Replacing generation with deterministic filtering removed two of the three
regressions, exactly as predicted: Hitler/Voss and Franco left the diff.


\section{Statistical Power of the Ablation Comparisons}\label{app:power}

One of four conditions reaches significance. That is the correct outcome to
report\index{McNemar} for four components tested on half-samples. It is
reported as such rather than rewritten into four confident claims.

A correction for multiple comparisons was not among the decisions fixed in
advance in \Cref{sec:pre-registration}. So it is applied here after the fact and
labelled as such. Sixteen paired tests are reported in this chapter: the eight
against the baselines in \Cref{sec:findings} and the eight ablation comparisons
above. Within the ablation family the one significant result survives correction
--- the planner effect on WebQSP is $p = 0.00592$ against a Bonferroni threshold
of $0.05/8 = 0.00625$. Being the smallest of the eight, it clears the first
step of Holm's procedure at that same threshold. Nothing else in the family is
near the bar under any correction. The next smallest is $p = 0.088$. The eight
baseline comparisons clear $0.05/8$ as well, the largest of them at $p =
0.0035$. Correction therefore changes no claim made here. That is the only
circumstance in which reporting it after the fact is worth anything.

The power limitation is arithmetic, not rhetorical. McNemar's test reads only
discordant pairs, and at $n \approx 200$ with a system near $0.5$--$0.75$
accuracy, a component that changes five questions in a hundred produces on the
order of ten discordant pairs --- not enough to reject at conventional levels.
The planner condition cleared the bar on WebQSP because its effect is large
($21$ against $6$), not because that comparison was better powered than the
others. The halves were adopted as the pre-registered response to the
benchmark's cost coming in above projection (\Cref{sec:ablation-conditions}).
Doubling them is the straightforward remedy and is named in
\Cref{sec:limitations-final}.

Three conditions are therefore reported as ``no detectable effect at this sample
size''. That phrase is not a hedge for ``no effect'': for the verifier the point
estimates are near zero \emph{and} the mechanism's contribution was never
expected to appear in accuracy. For backtracking and the $\alpha$-blend, by
contrast, the point estimates are consistently signed and the mechanism is
visibly active in the logs. Those are different epistemic situations. Collapsing
them into one word would lose the distinction.


\section{Building the Entity and Relation Indexes}\label{app:index-build}

\paragraph{Entity Embeddings}

Entity names were embedded with \textsf{all-MiniLM-L6-v2}, a 384-dimensional
sentence encoder~\cite{reimers2019sentence} chosen for CPU throughput at this
node count. Mediator nodes were skipped: embedding a string such as
\textsf{m.0y5k79m} produces noise. By \Cref{sec:graph-construction}, mediators
make up almost two thirds of the graph. So skipping them is also the
difference between a tractable and an intractable embedding job. Vectors were
written back into Neo4j and indexed with a cosine-similarity vector index.

\paragraph{Relation-Vocabulary Embeddings}

The scoring function of \Cref{sec:explorer} compares candidate relations against
the current sub-objective in embedding space. So the 7{,}058 relation names are
embedded once and held in memory rather than embedded per step.

Two details make this work. Relation names are \emph{verbalised} before
encoding: \textsf{people.\allowbreak person.\allowbreak place\_of\_birth}
becomes ``people person place of birth'', with consecutive duplicate tokens
removed. The encoder was trained on natural language. So the verbalised phrase
lands near a question like ``where was X born'' in embedding space, where the
raw dotted path does not. And relations are embedded with the \emph{same} model
as entities and sub-objectives, since vectors from different models occupy
unrelated spaces and their cosine similarity is meaningless.

The resulting artefact is a $7058 \times 384$ float32 matrix with a parallel
name list. Because vectors are L2-normalised at encoding time, similarity
against the whole vocabulary is a single matrix--vector product.

A functional check confirms the verbalisation is doing its job. Encoding the
probe ``where was the person born'' and ranking the vocabulary against it
returns \textsf{people.person.place\_of\_birth} first at cosine $0.709$ and
\textsf{location.location.people\_born\_here} --- the inverse relation ---
second at $0.675$, with the remaining top-five entries plausible near-misses
drawn from birth-related and fictional-character namespaces. A degenerate
verbaliser would have produced an unordered top five here.

\paragraph{Scope: Linking and Candidate Ranking, Never Ground
Truth}\label{sec:embedding-scope}

The role of embeddings in this thesis is deliberately narrow. They resolve
question phrases to graph entities, and they pre-rank relation candidates so the
language-model scorer can be applied to a short list rather than to thousands of
relations. They never establish that a fact holds. Every factual check ---
whether a triple exists, whether an entity is reachable, whether a claim is
supported --- is answered by symbolic graph queries. Nothing in the verification
layer of \Cref{sec:verification-layer} consults an embedding.


\section{Bulk Import into Neo4j}\label{app:bulk-import}

Import used the offline bulk importer of Neo4j Community 5.26.28, which operates
on a stopped database and is orders of magnitude faster than transactional
loading at this scale. Duplicate nodes and bad relationships were configured to
be skipped rather than to abort the import.

After import, a uniqueness constraint on \textsf{Entity.id} and a full-text
index on \textsf{Entity.name} were created. The full-text index is what the
resolver's lexical stage queries.


\section{Candidate Widths of the Agentic Baseline}\label{app:tog-widths}

\textbf{A second asymmetry is not in the baseline's favour, and it bears on how
\Cref{sec:tog-split} should be read.} Both systems call the same tools but
truncate the returned candidate sets at different widths: Think-on-Graph keeps
the first $40$ relations per entity and the first $20$ neighbours per relation,
the pruning widths of the reimplemented algorithm, against AGR's $300$ and $200$
(\Cref*{tab:budget}). Because \textsf{get\_relations} returns rows sorted by
\emph{descending} fanout (\Cref{sec:tool-api}), a binding cut discards the
low-fanout tail --- precisely the material \Cref{sec:tool-api} reports raising
AGR's own cap from $80$ to $300$ to stop discarding. Think-on-Graph therefore
runs at half of that $80$: half the cap this thesis found harmful for itself,
not half its own. Measured over the committed tool logs, the $40$-relation cut
binds on $31.6\%$ of the $1{,}651$ entities it expanded and the $20$-neighbour
cut on $32.8\%$ of its $7{,}992$ neighbour calls, against $1$ of AGR's $3{,}097$
expanded entities and $3.3\%$ of its neighbour calls. Beam width and depth
follow the original paper; the candidate widths behind them do not. They are a
property of this reimplementation. The consequence is a bound on one specific
claim. \Cref{sec:tog-split} concludes that the margin over Think-on-Graph is
affordability rather than search quality. That conclusion is safe in the
direction it is stated, since a narrower candidate set cannot explain why the
baseline runs \emph{out of calls}. But the residual gap on unclipped questions
is measured against a baseline searching a thinner pool. So it is a lower bound
on that baseline's attainable quality rather than an estimate of it. Equalising
the widths is the first item in \Cref{sec:future-baselines}.


\section{Verification Layer: A Worked Example}\label{app:verify-walkthrough}

\emph{Which university, founded in 1701, did actor James Franco attend?} Gold:
\textsf{University Yale}.

% TeX stops hyphenating a word at its first non-letter, so a dotted relation
% name can only ever break inside its first segment -- and it does, inserting a
% hyphen into the name. \allowbreak offers the dots as break points instead.
Navigation resolves \textsf{James Franco} as the anchor and expands along
\textsf{education.\allowbreak education.\allowbreak student} and
\textsf{people.\allowbreak person.\allowbreak education}. The plan's first
objective --- find the university founded in 1701 --- is never satisfied,
because the environment holds essentially no date literals
(\Cref{sec:unanswerable}). Three evaluator-triggered backtracks follow. The
run terminates on the depth budget with 18 model calls, 17 of them cache
replays.

The verifier drafts: \emph{``James Franco attended University Yale, which was
founded in 1701.''} Its entity list is \textsf{[University Yale]} --- copied
from the fact window, in the graph's spelling. It decomposes the draft into an
attendance claim and a founding claim.

The attendance claim resolves both endpoints in $\mu$. Its endpoint pair
is in the traversed adjacency set. It is supported. Four traversed triples
joining the pair are recorded as evidence. The founding claim names
\textsf{1701}. No traversed triple mentions it. So it cannot resolve, and
falls to the entailment check --- which correctly reports it unsupported,
because the fact is genuinely absent from the environment rather than merely
unretrieved. The verifier prepares a repair objective anchored on
\textsf{University Yale}. But the depth budget is already spent. So the router
takes \textsf{give\_up}, and the exploration never resumes. There was in any
case nothing to find.

The answerer now filters. \textsf{University Yale} appears in a supported claim.
So it is kept, verbatim. The rewrite call produces prose that asserts the
attendance and hedges the founding date. Its own opinion about entities is
discarded. Final answer entities: \textsf{[University Yale]} --- a hit.

The instructive part is the counterfactual. Before deterministic filtering, this
identical trace --- same draft, same claims, same verdicts --- ended with
\textsf{[Yale University]}. The rewrite call, asked to restate the answer,
restated it in the model's own spelling rather than the graph's. The answer
was scored wrong. Nothing about the verification was different. The claim-level
machinery had worked perfectly. The layer still emitted an entity string that
appears nowhere in the environment. That outcome is the precise thing it
exists to prevent.


\section{The Answerer's Exception Handler}\label{app:answerer-handler}

\Cref{sec:answerer} reports that the answerer's third path separates budget
exhaustion, which is expected, from a genuine error, which is not. What carries
that distinction is the recorded exception rather than the handler: the handler
catches \textsf{(BudgetExhausted, Exception)}, whose first arm is unreachable,
since the former subclasses the latter. It is the trace entry, which stores the
exception itself, that tells the two apart.

Whether it does so correctly is untested by anything in this work. Across all
$6{,}960$ committed runs no handler of this kind fired even once. The third path
is entered on a missing draft rather than on an exception. So it needs its own
check, and it gets the same answer: every one of the $3{,}690$ committed records
that ran this graph carries a non-empty draft.


\section{Structural-Check Route Frequencies}\label{app:route-frequencies}

Both routes are relation-agnostic. Since $r$ is free text with no correspondence
to a Freebase predicate, matching on it would reject true claims wholesale ---
``X played for Y'' is realised through a roster mediator carrying no predicate
named for playing. The design intent is a division of labour: the structural
routes provide high recall for ``some supporting edge exists.'' Semantic
precision is delegated to the entailment check. \Cref{sec:verify-failure-modes}
examines what that division actually buys and what it lets through.

On the development set the second route was consulted five times, across five
distinct questions --- four per cent of the 121 claims. All five returned
connected, and all five had a mediator path. In three of them there was no
direct edge at all. So a mediator-blind adjacency test would have rejected them.
Two were direct edges the traversal had simply never walked. Query latency
across those same five calls ranged from 1.9 to 75.7\,ms. With $n = 5$, that
range is a record of what happened rather than a latency distribution. The
route is therefore cheap and, on this evidence, correct --- but it is also rare.

Rare at test scale too. There the counts are large enough to carry weight.
Across the $800$ test questions the verifier decomposed $2{,}110$ claims over
$828$ firings, accepted $2{,}008$ of them and rejected $102$. The second route
was consulted on $100$ of those claims --- $4.7\%$, the same order as the
development set's four per cent --- and accepted $39$. Both figures are exact:
the first from the verifier's own trace entries, the second from the tool log,
which records one \textsf{verify\_connection} call per claim reaching that
branch because the code issues exactly one there.

What the record cannot say is how the remaining $1{,}969$ acceptances divide
between traversed adjacency and the entailment check. A claim routed to
entailment leaves no per-claim trace. The \textsf{n\_structural} counter that
\Cref{sec:output-contract} fences counts entailment-accepted claims alongside
structural ones. So it cannot separate them either. The traversed-adjacency
share is bounded above by $1{,}969$ and below by nothing at all. That is the
instrumentation gap of \Cref{sec:threats} in a third form. It is why this
paragraph reports the share of \emph{claims} that reached the second route,
which is measured, and not the share of \emph{acceptances} that came through
the first, which is not.


\section{Linking Misses and Gate Checks}\label{app:linking-misses}

The ten unresolved mentions of \Cref*{tab:linking} are not a random residue.
Every one is a bare machine identifier of the form \textsf{g.123852dg} appearing
as a CWQ topic entity --- an unnamed mediator node that the benchmark supplies
where a named entity was expected. Exact matching fails because no node carries
that string as a \emph{name}. The lexical stage then returns a meaningless
single-character match. The correct characterisation is that these questions are
anchored on nodes with no surface form, not that resolution is unreliable.

The strongest gate check re-verifies the Python-computed reachability against
the live database, confirming that nothing was lost between the CSV files and
the imported graph. A random sample of 400 questions flagged reachable, drawn
under a fixed seed, was re-tested with a shortest-path query in Neo4j. The pass
criterion
is strict --- these were already known to be reachable. So anything below 100\%
would indicate dropped edges. \textbf{All 400 verified. The failure set is
empty.} Together with the zero skipped import rows and the exact count
reconciliation, the environment is certified.


\section{Reconciling the Two Coverage
Measurements}\label{app:ceiling-reconciliation}

These ceilings are not the ones the validation gate reported. That difference
is worth naming rather than leaving for a reader to reconcile.
\Cref*{tab:coverage} measured the full test splits --- $1{,}628$ and $3{,}531$
questions --- and found $97.3\%$ and $99.7\%$ reachable. The figures above are
the same quantity measured over the $400$ questions actually evaluated, and they
come out at $97.0\%$ and $99.2\%$. Neither corrects the other: one bounds what
the environment could support across the whole benchmark, and the other bounds
what any system in \Cref{sec:results} could have scored. It is the second bound
that every accuracy in this thesis is read against. The gap between them is
sampling. It stays under a point on both datasets because the draw preserved
the strata proportions.


\section{The Test-Set Build That Read the Wrong
Split}\label{app:testset-nearmiss}

The glob was \textsf{train-*.parquet} where \textsf{test-*.parquet} was meant.
It came from the development-set miner, which reads train because it should.
The build then behaved like any other: eighteen shards resolved, no rows
skipped, pool sizes printed without complaint.

Two things were suggestive and neither forced the question. The pool sizes
matched the train splits, and eighteen shards resolved. Both are wrong only
relative to the test-split figures, and reading them that way means already
suspecting the split.

What forced it was an assertion rather than an observation. The
development-overlap guard reported $32$ overlapping questions on WebQSP and
$45$ on CWQ against a required zero: the development set is mined from the
train and validation splits (\Cref{sec:dev-construction}), so a pool that
overlaps it at all is not a test pool. A threshold fires whether or not anyone
is reading for it, which is why it settled what the two suggestive signals had
not.


\section{When the Retry Route Can Fire}\label{app:retry-budget}

The retry route requires verify, depth, \emph{and} time budget. The depth
condition was added after a development trace showed a question completing with
depth 5 against a maximum of 4: the retry edge re-enters the explorer without
passing through the post-evaluation router, which is where the depth check
lives. So each retry silently granted one expansion beyond budget. Since
\Cref{sec:budgets} claims budgets are hard guarantees, the hole was closed in
the router rather than documented as an exception.

That condition turns out to be decisive on the development set, where
\textbf{the route never executed}. On all three development questions that
produced an unsupported claim, the depth budget was already exhausted when the
verifier ran. So the router took \textsf{give\_up} in every case, and no repair
exploration occurred. That pattern is not a coincidence of small numbers so
much as a structural correlation: the verifier is reached either because the
plan completed or because a budget ran out. All three firings came from
the fourteen questions that exhausted depth, none from the sixty-six that
finished with search budget in hand.

\textbf{That correlation does not survive the move to test, and the route
executes there.} Across the $800$ test questions the verifier asks for a
repair on $52$ --- $16$ on WebQSP and $36$ on CWQ. Exploration actually
resumes on $28$ of them, $10$ and $18$ respectively. So the development
observation was about a sample of three. Generalising it would have been
wrong: the depth condition blocks the retry about half the time rather than
always. What this chapter cannot answer is whether the route \emph{helps}.
The census answers that in the direction the design did not intend ---
\Cref{sec:verifier-errors} reads four of the five wrongly-rejected specimens
as retries that ran and lost the answer. That is where the recall cost
\Cref{sec:verification-tradeoff} attributes to this layer is actually paid.
What the development set establishes is therefore narrower than it looked:
that the route contributed nothing \emph{there}, not that nothing in this
thesis depends on it.

One instrumentation consequence follows, and matters when reading run records.
The verifier increments \textsf{verify\_iters} and appends the repair objective
\emph{before} the router decides. So a record showing \textsf{verify\_iters:\,1}
records the verifier's intent to repair, not an executed repair. Whether
exploration actually resumed is visible only in the node sequence of the trace.


\section{The Resolution Filter's Head--Tail
Asymmetry}\label{app:filter-asymmetry}

\textbf{The table is built from tail names only, and the verifier's is not.} A
traversed triple contributes its tail to this lookup and not its head. So an
entity the search \emph{departed from} --- a topic entity, or any intermediate
the agent has since moved past --- cannot be resolved as an answer by the
evaluator, however the model phrases it. The structural check of
\Cref{sec:structural-check} takes the opposite convention: it collapses
traversed triples into undirected endpoint pairs. So heads and tails are
interchangeable there.

The asymmetry is deliberate in the evaluator and worth naming because it is not
obviously so. A topic entity is exactly what an echoing model reaches for
(\Cref{sec:echo}). Refusing to resolve it is a cheap structural defence against
the most common degenerate answer. The cost is that a question whose gold
answer genuinely \emph{is} a topic entity, or an intermediate the plan has
passed through, cannot be answered through this path at all. Those cases
exist --- \Cref{sec:echo} counts $13$ \textsf{echo} failures. The filter is
what prevents them from being far more numerous. But the two layers
disagreeing about what counts as a reachable entity is a design seam, not a
proved optimum. No experiment in this thesis isolates it.


\section{Why verify\_triple Is Not Used}\label{app:verify-triple}

\Cref*{tab:tools} lists the operations. There are five, not the four originally
planned. Only four of the five are reachable from a node in the final design.
Both facts are consequences of the same episode and are recorded rather than
tidied away.

\textsf{verify\_triple} was built first, as the claim checker: given a claim
$\langle h, r, t\rangle$ it asked whether that relation held between those
entities. It did not survive contact with generated text. Claim decomposition
produces relations phrased in English --- ``played for'', ``is the capital
of''. Freebase predicates are neither phrased that way nor reliably one edge
long. So requiring the relation to match rejected claims that were true.
\textsf{verify\_connection}\index{verify connection@\textsf{verify\_connection}}
was added in response. It drops the relation entirely and asks only about
adjacency. This version answered the question the verifier actually needed
answered, and it superseded rather than complemented its predecessor.
\Cref{sec:structural-check} describes the two routes the verifier ended up
with. \textsf{verify\_triple} is not one of them.

The operation is left in the API because it is exercised by the tool
integration tests and remains the natural primitive for a relation-aware
check. \Cref{sec:future} proposes exactly that check. But nothing in the
running loop calls it. The table says so.

\endinput
